\documentclass[a4paper,11pt]{article}

% --- Pakker for matematikk og font ---
\usepackage[english]{babel}
\usepackage{amsmath, amsthm, amssymb}
\usepackage{mathtools} 
\usepackage{geometry} 
\usepackage{enumitem}
\geometry{margin=2cm}

\setlength{\parindent}{0pt}
\setlength{\parskip}{0.5em} 

% make amsthm use a filed in black square
\renewcommand{\qedsymbol}{$\blacksquare$}

% Mengder
\newcommand{\N}{\mathbb{N}}
\newcommand{\Z}{\mathbb{Z}}
\newcommand{\Q}{\mathbb{Q}}
\newcommand{\R}{\mathbb{R}}

% Topologi
\newcommand{\T}{\mathcal{T}} 
\newcommand{\B}{\mathcal{B}}
\newcommand{\inv}[1]{#1^{-1}} 
\newcommand{\interior}[1]{#1^\circ} 
\newcommand{\closure}[1]{\overline{#1}} 
\newcommand{\boundary}{\partial} 

% Greske bokstaver
\newcommand{\eps}{\varepsilon}

% --- Document Info ---
\title{MAT3500 - Topology: \\ Mandatory assignment}
\author{Vegard Nyquist Borgeraas}
\date{Autumn 2026}

\begin{document}

\maketitle

\section*{Problem 1}
Let $(X, d)  $ be a metric space.

\begin{enumerate}[label=\textbf{(\alph*)}]
    \item Show that $d \colon X \times X \to \R$ is continuous, where $X \times X $ has the product topology.

    \begin{proof}
    Let $B := (a,b)$ be a basis element in $\R$, and $(x,y) \in \inv{d}(B)$.
    
    Then let $\eps := \min\left(\frac{b-d(x,y)}{2}, \frac{d(x,y)-a}{2}\right)$.
    
    Let $U := B_d(x,\eps) \times B_d(y,\eps)$ be an open neighbourhood around $(x,y)$.
    Let $(x',y') \in U$.
    
    Then by applying the triangle inequality twice, we get:
    \[ d(x',y') \leq d(x,y) + d(x',x) + d(y,y') < d(x,y) + 2\eps \leq b\]
    And by applying the inverse triangle inequality twice we get:
    \[ d(x',y') \geq d(x,y) - d(x',x) - d(y,y') > d(x,y) - 2\eps \geq a\]
    
    
    Therefore $(x',y') \in \inv{d}(B)$, which implies $U \subseteq \inv{d}(B)$. Thus, $\inv{d}(B))$ is open.
    
    Then we have that $d$ is continuous.
    \end{proof}

    \item Show that for $x,y \in X$, subset $A:=\{z \in X : d(x,z) = d(y,z)\}$ is closed in $X$.

    \begin{proof}
    Let $f \colon X \to \R$ be defined by $z \mapsto d(x,z)-d(y,z)$, then we have that $A = \inv{f}(\{0\})$. Since $d$ is continuous it is also continuous in each variable separately, then $f$ is continuous. And since $\{0\}$ is closed, then so is $A$.
    \end{proof}
\end{enumerate}

\newpage

\section*{Problem 2}

    Let $S := \N \cup \{\infty\}$ and $X$ be a topological space. Find a topology on $S$ for which a function $f \colon S \to X$ is continuous if and only if the sequence $f(n)$ converges to $f(\infty)$ in $X$.

\begin{proof}
    Let $S$ have the topology generated by the subbasis $(\{n\})_{n \in \N} \cup (\{n \in S: n \geq N \})_{N \in \N}$.

    First, assume that $f$ is continuous. 
    Let $U \subseteq X$ be an open neighborhood of $f(\infty)$. 
    Since $f$ is continuous, there exists a basis element $B$ such that $\infty \in B$ and $f(B) \subseteq U$. 
    Since $\infty \in B$, $B$ is a finite intersection of subbasis elements of the form $\{n \in S : n \geq N_i\}$. 
    Thus, $B = \{n \in S : n \geq N\}$ where $N = \max(N_i)$.
    Therefore, $f(n) \in U$ for all $n \geq N$.
    Then $f(n) \to f(\infty)$. 

    Conversely, assume that $f(n) \to f(\infty)$ as $n \to \infty$. 
    Let $U$ be open in $X$, and let $s \in \inv{f}(U)$. 
    If $s \neq \infty$, then the singleton $\{s\}$ is open, and $s \in \{s\} \subseteq \inv{f}(U)$. 
    If $s = \infty$, by our assumption there exists an $N \in \N$ such that $f(n) \in U$ for all $n \geq N$. 
    Then the open set $\{n \in S: n \geq N \} \subseteq \inv{f}(U)$. 
    Thus, $\inv{f}(U)$ is open, and $f$ is continuous. 

\end{proof}


\section*{Problem 3}
Let $X, Y$ be topological spaces, $f, g \colon X \to Y$ two continuous functions, and assume that $f|_A = g|_A$ for a dense subset $A \subseteq X$.

\begin{enumerate}[label=\textbf{(\alph*)}]
    \item Show that if $Y$ is Hausdorff, then $f = g$.
    \begin{proof}
        Assume for contradiction that $f \neq g$.
        Then there exists an $x \in X$ such that $f(x) \neq g(x)$. \\
        Since $Y$ is Hausdorff, there exist open neighbourhoods $V_f$ and $V_g$ of $f(x)$ and $g(x)$ respectively, such that $V_f \cap V_g = \emptyset$.
        Since $f$ and $g$ are continuous, there exist open neighbourhoods $U_f$ and $U_g$ of $x$ such that $f(U_f) \subseteq V_f$ and $g(U_g) \subseteq V_g$.
        Define $U := U_f \cap U_g$. Then $U$ is an open neighbourhood of $x$. 
        Then we have $f(U) \cap g(U) \subseteq V_f \cap V_g$. 
        Since $A$ is dense, there exists a $y \in U \cap A$. Since $y \in A$, we have $f(y) = g(y)$. 
        Then $f(y) \in f(U) \cap g(U)$, which implies $V_f \cap V_g \neq \emptyset$. This is a contradiction. Thus, $f=g$.
    \end{proof}

    \item Find an example which shows that this is not always true if $Y$ is not Hausdorff.
    \begin{proof}
        Let $X = Y = \{a,b\}$ be equipped with the trivial topology, and let $A=\{a\}$.
        Then $Y$ is not Hausdorff and $A$ is dense in $X$.
        Let $f(x)=a$ and $g(x)=x$. Then $f|_A = g|_A$. \\
        However, $f \neq g$.
    \end{proof}
\end{enumerate}

\section*{Problem 4}
Consider the following collection of subsets of $\R$:
\[\B = \{[a,\infty) \subseteq \R \colon a \in \R\}\]

\begin{enumerate}[label=\textbf{(\alph*)}]
    \item Show that $\B$ forms a basis for a topology $\T'$ on $\R$.
    \begin{proof}
        First, let $x \in \R$. Then $x \in [x,\infty) \in \B$.

        Second, let $[a_1,\infty), [a_2,\infty) \in \B$, and define $a := \max(a_1,a_2)$.
        The intersection of these two sets is exactly $[a_1,\infty) \cap [a_2,\infty) = [a,\infty)$.
        Therefore, for any $x \in [a_1,\infty) \cap [a_2,\infty)$, we have $x \in [a,\infty) \in \B$, and clearly $[a,\infty) \subseteq [a_1,\infty) \cap [a_2,\infty)$.
        Thus, $\B$ forms a basis.
    \end{proof}

\item Show that every continuous function $f\colon (\R, \T') \to (\R, \T)$ is constant, where $\T$ denotes the usual topology on $\R$.
\begin{proof}
        Assume for contradiction that there exist $x,y \in \R$ such that $f(x) \neq f(y)$.
        Then the singletons $\{f(x)\}$ and $\{f(y)\}$ are disjoint and closed in $(\R, \T)$. \\
        Since $f$ is continuous, the preimage of a closed set is closed. 
        All open sets in $(\R, \T')$ are unions of basis elements, so they are of the form $\emptyset$, $\R$, $(a,\infty)$, or $[a,\infty)$ for some $a \in \R$. 
        Thus, the closed sets in $(\R, \T')$ are their complements, which are of the form $\emptyset$, $\R$, $(-\infty,b]$, or $(-\infty,b)$ for some $b \in \R$.
        
        Since $x \in \inv{f}(\{f(x)\})$, we have $\inv{f}(\{f(x)\}) \neq \emptyset$. The same holds for $\inv{f}(\{f(y)\})$.
        Notice that any non-empty closed set in $(\R, \T')$ must contain a set of the form $(-\infty, b)$ for some $b \in \R$.
        Therefore, there exist $b_x, b_y \in \R$ such that $(-\infty, b_x) \subseteq \inv{f}(\{f(x)\})$ and $(-\infty, b_y) \subseteq \inv{f}(\{f(y)\})$.
        Let $b = \min(b_x, b_y)$, then $(-\infty, c) = (-\infty, b_x) \cap (-\infty, b_y)$. \\
        Then we have:
        \[ \emptyset \neq (-\infty, b) \subseteq \inv{f}(\{f(x)\}) \cap \inv{f}(\{f(y)\}) = \inv{f}(\{f(x)\} \cap \{f(y)\}) = \inv{f}(\emptyset) = \emptyset \]
        This is a contradiction. Thus, $f$ must be constant.
    \end{proof}
\end{enumerate}

\section*{Problem 5}
Let $A = S^1 \cap \{(x, y) \in \R^2 \mid y \geq 0\}$, $B = \{(x, y) \in \R^2 \mid x \in [-1, 1], y = 0\}$, and $X = A \cup B$ be their union, with the subspace topology as a subspace of $\R^2$. Show that $X$ is homeomorphic to $S^1$.
\begin{proof}
\end{proof}

\section*{Problem 6}
Define an equivalence relation on $\R$ as follows:
\[ x \sim y \iff x - y \in \Z. \]
Write $\R/\Z$ for the quotient $\R/{\sim}$. Show that $\R/\Z$ is homeomorphic to $S^1$.

\section*{Problem 7}
Let $X = \prod_{n\in\N}\{-1, 1\} = \{-1, 1\}^\N$ with the product topology (where $\{-1, 1\}$ has the discrete topology).

\begin{enumerate}[label=\textbf{(\alph*)}]
    \item For $0 < q < 1$, show that $f \colon X \to \R$,
    \[ f((x_n)_{n\in\N}) = \sum_{n=1}^{\infty} x_n q^n \]
    is continuous.
    
    \item For $0 < q < \frac{1}{2}$, show that $f$ is injective.
    
    \item For $0 < q < \frac{1}{2}$, show that $f$ is actually an embedding, by showing that $f \colon X \to f(X)$ is open.
\end{enumerate}

\end{document}